\section{Funzioni dell'orchestratore e delle chiamate ai modelli}
\label{sec:orchestrator-implementation}
Questa sezione riproduce estratti del codice corrente. Le funzioni locali sono
definite dentro \texttt{run\_global\_goal\_orchestrator} e usano il suo stato tramite
closure: non sono procedure autonome da eseguire separatamente.
La funzione principale mantiene il controllo della sequenza descritta nella
Sezione~\ref{sec:goal-orchestrator}. Non sono richieste conversazioni tra LLM
per decidere quale stadio invocare.

\clearpage
\begin{lstlisting}[style=llmprompt,language=Python,caption={Definizione integrale del punto di ingresso dell'orchestratore: configurazione, fasi, aggiornamento e arresto.},label={lst:orchestrator-main-full}]
def run_global_goal_orchestrator(
    *, project_description: str, initial_high_level_goals: HighLevelGoals,
    initial_low_level_goals: LowLevelGoals, mode=ShotPromptingMode.ZERO_SHOT,
    evaluator_ablation=False, evaluation_output_directory=None,
    max_iterations=MAX_ITERATIONS, actors: Actors | None = None,
    verbose: bool = False, output_file_prefix: str | None = None,
    evaluator_few_shot_examples: int | None = None,
    evaluator_temperature: float | None = None,
    parent_threshold: float | None = None, duplicate_threshold: float | None = None,
    on_state=None,
) -> GlobalGoalCycleResult:
    """Coordinate reconstruction, documentary votes and approved top-down repairs."""
    # [01] Verifica le condizioni della pipeline top-down riusata.
    if mode != ShotPromptingMode.ZERO_SHOT or evaluator_ablation:
        raise ValueError("Bottom-up repairs require original zero-shot generation and an enabled critic")
    if type(max_iterations) is not int or not 1 <= max_iterations <= MAX_ITERATIONS:
        raise ValueError("max_iterations must be between 1 and 25")
    evaluator_few_shot_examples = (CRITIC1_FEW_SHOT_EXAMPLE_COUNT if evaluator_few_shot_examples is None
                                   else evaluator_few_shot_examples)
    evaluator_examples(evaluator_few_shot_examples)
    evaluator_temperature = BOTTOM_UP_EVALUATOR_TEMPERATURE if evaluator_temperature is None else evaluator_temperature
    if not 0 <= evaluator_temperature <= 2:
        raise ValueError("Evaluator temperature must be between 0 and 2")
    similarity_options = {"parent_threshold": parent_threshold, "duplicate_threshold": duplicate_threshold}
    active_similarity = similarity_configuration(**similarity_options)
    file_prefix = ""
    if output_file_prefix is not None:
        safe_name = re.sub(r"[^\w.-]+", "_", output_file_prefix).strip("._")
        if not safe_name:
            raise ValueError("Output file prefix must contain a project name")
        file_prefix = safe_name + "_"
    # [02] Copia la baseline: le revisioni lavorano sullo stato corrente.
    current_hlgs = _coerce_model(initial_high_level_goals, HighLevelGoals, context="initial HLGs").model_copy(deep=True)
    current_llgs = _coerce_model(initial_low_level_goals, LowLevelGoals, context="initial LLGs").model_copy(deep=True)
    known_actors = _merge_actors(actors or Actors(actors=[]), _actors_from_goals(current_hlgs))
    traces, warnings = [], []
    initial_adjustments = []
    # [03] Riconcilia i riferimenti prima del raggruppamento dei figli.
    current_llgs = _reconcile_parent_references(
        current_hlgs, current_llgs, strict=False, adjustments=initial_adjustments, warnings=warnings)
    # [04] Registra policy e parametri per ricostruire questa esecuzione.
    config = {"pipeline_policy_version": PIPELINE_POLICY_VERSION,
              "orchestrator_module": "src.bottom_up.goal_orchestrator",
              "execution_structure": "functions_over_top_down_models",
              "reconstruction_history_policy": "fresh_messages_per_call_and_retry_no_prior_responses",
              "output_file_prefix": file_prefix,
              "evaluator_models": list(BOTTOM_UP_EVALUATOR_MODELS),
              "evaluator_temperature": evaluator_temperature,
              "evaluator_few_shot_examples": evaluator_few_shot_examples,
              "evaluator_response_format": "json_schema",
              "evaluator_endpoint": "/v1/chat/completions",
              "evaluator_retry_policy": "whole_global_response_per_model",
              "evaluator_attempts_per_iteration": BOTTOM_UP_EVALUATOR_VOTE_ATTEMPTS,
              "context_recovery_policy": "reload_affected_model_double_context",
              "similarity_technical_attempts": SIMILARITY_TECHNICAL_ATTEMPTS,
              "review_payload_format": "explicit_prime_original_candidate_pairs_v1",
              "response_length_policy": "rewrite_feedback_500_character_target_no_length_rejection",
              "similarity": active_similarity,
              "hierarchy_input_policy": "original_top_down_hlg_and_llg_collections",
              "goal_counter_scope": "position_within_iteration_snapshot",
              "max_global_iterations": max_iterations,
              "repair_policy": "original_top_down_start_at_hlg_then_llgs",
              "repair_feedback_policy": "join_agreeing_feedback_critique_only",
              "repair_acceptance": "hlg_critic_approval_then_llg_structural_validation_atomic_commit",
              "parent_reference_policy": "exact_reference_first_then_unique_exact_description",
              "initial_parent_reference_reconciliations": initial_adjustments,
              "repair_reflection_max_attempts": MAX_ATTEMPTS,
              "repair_quality_threshold": QUALITY_THRESHOLD,
              "repair_prompting_mode": "ZERO_SHOT",
              "repair_critic_enabled": True,
              "hlg_collection_repair_policy": "one_batch_pass_per_iteration_until_global_stop",
              "removal_schedule": "all_approved_removals_before_batch_rewrites",
              "generator_intervention_policy": "full_collection_generation_no_bottom_up_freeze",
              "action_schedule": "all_approved_removals_then_all_surviving_rewrite_suggestions",
              "output_verbosity": "detailed" if verbose else "iteration_summary"}

    # [05] Pubblica uno stato finale oppure un checkpoint ancora parziale.
    def publish(pending=None, result=None):
        if result is not None:
            state = result.model_dump(mode="json", exclude_defaults=True)
        else:
            state = dict(schema_version=RESULT_SCHEMA_VERSION, configuration=config, completed_iterations=len(traces),
                         warnings=list(warnings),
                         final_high_level_goals=current_hlgs.model_dump(mode="json"),
                         final_low_level_goals=current_llgs.model_dump(mode="json"),
                         iterations=[trace.model_dump(mode="json", exclude_defaults=True)
                                     for trace in traces + ([pending] if pending is not None else [])])
        if evaluation_output_directory is not None:
            atomic_write(Path(evaluation_output_directory) / f"{file_prefix}results.json", compact_cycle(state))
        if on_state is not None:
            on_state(state)

    # [06] Costruisce il risultato con il motivo di arresto.
    def finish(reason, converged):
        result = GlobalGoalCycleResult(
            schema_version=RESULT_SCHEMA_VERSION,
            converged=converged, stop_reason=reason, completed_iterations=len(traces),
            final_high_level_goals=current_hlgs, final_low_level_goals=current_llgs,
            iterations=traces, warnings=warnings, configuration=config)
        publish(result=result)
        return result

    publish()

    # [07] Il limite globale resta fisso; aumenta solo il contatore.
    for iteration in range(1, max_iterations + 1):
        print(f"[bottom-up] iteration {iteration}/{max_iterations} | full snapshot review", flush=True)

        def run_stage(name, function, *args, **kwargs):
            if verbose:
                return function(*args, **kwargs)
            with redirect_stdout(io.StringIO()), redirect_stderr(io.StringIO()):
                return function(*args, **kwargs)

        # [08] Gli ID di ramo appartengono a questo snapshot.
        branches = indexed_high_level_goals(current_hlgs)
        children_by_hlg = {bid: low_level_goals_for_hlg(parent, current_llgs)
                           for bid, parent in branches.items()}
        branch_warnings = [f"HLG_WITHOUT_CHILDREN: branch {bid} ({branches[bid].name})"
                           for bid, children in children_by_hlg.items() if not children]
        if verbose:
            for message in branch_warnings:
                print(f"[bottom-up] {message}", flush=True)
        # [09] Qwen ricostruisce HLG prime soltanto per gruppi non vuoti.
        reconstructions = {k: _coerce_model(v, BottomUpHighLevelGoal, context="reconstruction")
                           for k, v in run_stage("reconstruction", reconstruct_all_branches,
                               {bid: children for bid, children in children_by_hlg.items() if children}).items()}
        # [10] MPNet calcola relazioni e candidati senza un giudizio LLM.
        analysis = run_stage("similarity", analyze_similarities, branches, reconstructions, **similarity_options)
        config["similarity"] = analysis["configuration"]
        cosine_payload = build_cosine_payload(branches, reconstructions, analysis, iteration=iteration)
        if verbose:
            print(json.dumps(cosine_payload, ensure_ascii=False, indent=2), flush=True)
        global_votes = []

        # [11] Salva ogni risposta di modello prima degli aggiornamenti.
        def save_global_vote(record):
            global_votes.append(record)
            checkpoint(0, len(branches), {})

        def checkpoint(done, total, evaluations):
            # [12] Costruisce un checkpoint parziale con dati e voti correnti.
            trace = GlobalGoalCycleIteration(
                iteration=iteration, reconstructions=reconstructions, evaluations=evaluations,
                high_level_goals=current_hlgs, low_level_goals=current_llgs,
                evaluator_responses=list(global_votes), similarity_analysis=analysis,
                input_hierarchy={"high_level_goals": current_hlgs.model_dump(mode="json"),
                                 "low_level_goals": current_llgs.model_dump(mode="json")},
                checkpoint_complete=False, completed_branch_count=done, total_branch_count=total)
            publish(pending=trace)

        checkpoint(0, len(branches), {})
        # [13] I tre evaluator ricevono lo stesso snapshot e contesto.
        evaluations = run_stage("documentary_review", evaluate_all_branches,
            project_description, branches, reconstructions, actors=known_actors,
            on_global_vote=save_global_vote,
            on_evaluation=checkpoint, similarity_analysis=analysis, cosine_payload=cosine_payload,
            few_shot_example_count=evaluator_few_shot_examples, evaluator_temperature=evaluator_temperature)
        evaluations = {k: _coerce_model(v, GlobalGoalEvaluationResult, context="global review")
                       for k, v in evaluations.items()}
        start_hlgs, start_llgs = current_hlgs, current_llgs
        # [14] Applica le sole proposte approvate e tenta la riparazione.
        update = apply_review_decisions(
            project_description=project_description, known_actors=known_actors,
            current_hlgs=current_hlgs, current_llgs=current_llgs,
            branches=branches, evaluations=evaluations,
            repair=lambda **kwargs: run_stage(
                "top_down_repair", _repair_collections_with_top_down, **kwargs))
        current_hlgs, current_llgs = update["high_level_goals"], update["low_level_goals"]
        evaluations, outcomes = update["evaluations"], update["outcomes"]
        changed, confirmed = update["changed"], update["confirmed"]
        top_down_reflections, iteration_warnings = update["reflections"], update["warnings"]
        iteration_warnings.extend(branch_warnings)
        if verbose:
            for bid, outcome in outcomes.items():
                print(f"[bottom-up] applied | branch={bid} | {outcome}", flush=True)
        context_blocked = not changed and any(
            result.evaluation_status == "CONTEXT_LENGTH_EXCEEDED" for result in evaluations.values()
        )
        if context_blocked:
            iteration_warnings.append("EVALUATOR_CONTEXT_LENGTH_EXCEEDED: iteration attempt budget exhausted; remaining valid actions applied; review continues next iteration")
        if branches and all(analysis["branches"][b]["status"] == "ERROR" for b in branches):
            iteration_warnings.append("SIMILARITY_UNAVAILABLE: iteration attempt budget exhausted; review continues next iteration")
        trace = GlobalGoalCycleIteration(
            iteration=iteration, reconstructions=reconstructions, evaluations=evaluations,
            top_down_reflections=top_down_reflections,
            evaluator_responses=global_votes, similarity_analysis=analysis,
            high_level_goals=current_hlgs, low_level_goals=current_llgs,
            input_hierarchy={"high_level_goals": start_hlgs.model_dump(mode="json"),
                             "low_level_goals": start_llgs.model_dump(mode="json")},
            structural_action_outcomes=outcomes,
            warnings=iteration_warnings, checkpoint_complete=True,
            completed_branch_count=len(evaluations), total_branch_count=len(branches))
        # [15] Archivia la traccia completa solo dopo gli aggiornamenti della gerarchia.
        traces.append(trace)
        warnings.extend(iteration_warnings)
        publish()
        print(f"[bottom-up] iteration {iteration} complete | branches={len(branches)} | "
              f"changed={changed} | decisions={dict(Counter(e.decision.value if e.decision else 'INCONCLUSIVE' for e in evaluations.values()))}",
              flush=True)
        if evaluation_output_directory is not None:
            print(f"  updated: {file_prefix}results.json", flush=True)
        # [16] Arresta per conferma solo senza cambiamenti e con tutti i rami confermati.
        if not changed and len(confirmed) == len(branches):
            return finish(GlobalGoalCycleStopReason.ALL_BRANCHES_CONFIRMED, True)
    # [17] Conserva il risultato anche se il limite viene raggiunto senza convergenza.
    return finish(GlobalGoalCycleStopReason.MAX_ITERATIONS_REACHED, False)
\end{lstlisting}

\begin{lstlisting}[style=llmprompt,language=Python,caption={Applicazione delle decisioni: rimozioni approvate, raccolta delle critiche e riparazione batch.},label={lst:orchestrator-apply-full}]
def apply_review_decisions(
    *, project_description: str, known_actors: Actors,
    current_hlgs: HighLevelGoals, current_llgs: LowLevelGoals,
    branches, evaluations, repair,
):
    """Remove approved branches, then send all rewrite feedback to top-down.

    Inputs are the existing goal models and evaluator outputs. Return updated
    collections and trace data; do not constrain the generator's intervention.
    """
    # [01] Lavora sulle collezioni correnti indicizzate per ramo.
    next_goals = dict(branches)
    removed_parents = []
    outcomes, iteration_warnings = {}, []
    changed = False
    # [02] Protegge i rami indicati come superstiti da rimozioni incompatibili.
    protected = {e.surviving_branch_id for e in evaluations.values()
                 if e.surviving_branch_id}
    for branch in branches:
        if (evaluations[branch].decision in {HighLevelGoalDecision.REMOVE, HighLevelGoalDecision.DUPLICATE_REMOVAL}
                and branch in protected):
            evaluations[branch] = _mark_unresolved(evaluations[branch], "Conflicting survivor removal.")
    for branch in branches:
        result = evaluations[branch]
        if result.decision not in {HighLevelGoalDecision.REMOVE, HighLevelGoalDecision.DUPLICATE_REMOVAL}:
            continue
        # [03] Rimuove il padre; i figli associati sono filtrati subito dopo.
        next_goals.pop(branch)
        removed_parents.append(branches[branch])
        outcomes[branch] = "REMOVED_WITH_CHILDREN"
        changed = True

    # Apply all approved removals before repairing the surviving collection.
    current_hlgs = HighLevelGoals(goals=list(next_goals.values()))
    current_llgs = LowLevelGoals(low_level_goals=[
        child for child in current_llgs.low_level_goals
        if not any(child.high_level_associated.model_dump(mode="json") == parent.model_dump(mode="json")
                   for parent in removed_parents)
        or any(child.high_level_associated.model_dump(mode="json") == parent.model_dump(mode="json")
               for parent in next_goals.values())])
    # [04] Raccoglie le critiche dei rami superstiti da riscrivere.
    hlg_requests, requested_ids = [], []
    top_down_reflections = {}
    for branch in branches:
        bid = branch
        if bid not in next_goals:
            continue
        result = evaluations[bid]
        rewrite = result.decision == HighLevelGoalDecision.REWRITE
        if not rewrite:
            continue
        instruction = result.feedback.critique if result.feedback else None
        if not instruction or not instruction.strip():
            evaluations[bid] = _mark_unresolved(result, "Missing repair instruction")
            outcomes[bid] = "REJECTED_STATE_PRESERVED"
            iteration_warnings.append(f"{bid}: Missing repair instruction")
            continue
        hlg_requests.append(instruction)
        requested_ids.append(bid)
    handoff_completed = False
    # [05] Una sola riparazione batch riceve tutti i feedback raccolti.
    if requested_ids:
        try:
            candidate_hlgs, candidate_llgs = repair(
                description=project_description, actors=known_actors,
                high_level_goals=current_hlgs, low_level_goals=current_llgs,
                hlg_requests=hlg_requests,
                records=top_down_reflections)
            # [06] Confronta i contenuti: handoff riuscito non implica cambiamento.
            state_changed = candidate_hlgs != current_hlgs or candidate_llgs != current_llgs
            current_hlgs, current_llgs = candidate_hlgs, candidate_llgs
            changed = changed or state_changed
            handoff_completed = True
            for bid in requested_ids:
                outcomes[bid] = "HLG_LLG_HANDOFF_COMPLETED"
        # [07] Se la riparazione fallisce conserva lo stato dopo le rimozioni.
        except Exception as exc:
            message = f"REPAIR_REJECTED: {type(exc).__name__}: {exc}"
            for bid in requested_ids:
                outcomes[bid] = "REJECTED_STATE_PRESERVED"
            iteration_warnings.append(message)
    
    # [08] Dopo un handoff la nuova gerarchia richiede una nuova verifica.
    confirmed = [] if handoff_completed else [bid for bid, e in evaluations.items()
                 if bid in next_goals and e.decision == HighLevelGoalDecision.KEEP]
    # [09] Assegna un esito anche ai rami senza interventi strutturali.
    for bid, result in evaluations.items():
        outcomes.setdefault(bid, result.decision.value if result.decision else "INCONCLUSIVE")

    return dict(high_level_goals=current_hlgs, low_level_goals=current_llgs,
                evaluations=evaluations, outcomes=outcomes, changed=changed,
                reflections=top_down_reflections, warnings=iteration_warnings,
                confirmed=confirmed)
\end{lstlisting}

\begin{lstlisting}[style=llmprompt,language=Python,caption={Validatore operativo corrente delle singole decisioni: i controlli strutturali sono distinti dalle condizioni documentali del prompt.},label={lst:documentary-decision-validator}]
def _validate_decision(item, branch, analysis):
    vote = DocumentaryDecision.model_validate(item)
    if vote.branch_id != branch:
        raise ValueError("Wrong branch_id")
    record = analysis["branches"][branch]
    if record["status"] != "VALID":
        raise ValueError("Similarity unavailable; branch must remain inconclusive")
    if vote.decision == "DUPLICATE_REMOVAL":
        if vote.surviving_branch_id not in record["duplicate_candidates"]:
            raise ValueError("Survivor must be a direct same-actor text-review candidate")
    elif vote.decision == "REWRITE_ORIGINAL_HIGH_LEVEL_GOAL":
        if not record["reconstructable"] or record["relationship"] not in {"EQUIVALENT", "NOT_EQUIVALENT"}:
            raise ValueError("REWRITE requires a reconstructable HLG prime and a valid numerical relationship")
        if record.get("original_hlg_comparisons_complete") is False:
            raise ValueError("Cannot choose REWRITE before completing available reconstructed-HLG comparisons")
    elif vote.decision == "REMOVE_ORIGINAL_HIGH_LEVEL_GOAL":
        if record["relationship"] != "EQUIVALENT":
            raise ValueError("REMOVE requires an EQUIVALENT parent relationship")
    return vote
\end{lstlisting}

\clearpage
\begin{lstlisting}[style=llmprompt,language=Python,caption={Wrapper che esegue una funzione di fase e controlla soltanto la visibilità dei log.},label={lst:orchestrator-run-stage}]
def run_stage(name, function, *args, **kwargs):
    if verbose:
        return function(*args, **kwargs)
    with redirect_stdout(io.StringIO()), redirect_stderr(io.StringIO()):
        return function(*args, **kwargs)
\end{lstlisting}

\begin{lstlisting}[style=llmprompt,language=Python,caption={Callback e salvataggio dello stato: i checkpoint parziali restano distinti dalle iterazioni complete.},label={lst:orchestrator-state-callbacks}]
def save_global_vote(record):
    global_votes.append(record)
    checkpoint(0, len(branches), {})

def checkpoint(done, total, evaluations):
    trace = GlobalGoalCycleIteration(
        iteration=iteration, reconstructions=reconstructions, evaluations=evaluations,
        high_level_goals=current_hlgs, low_level_goals=current_llgs,
        evaluator_responses=list(global_votes), similarity_analysis=analysis,
        input_hierarchy={"high_level_goals": current_hlgs.model_dump(mode="json"),
                         "low_level_goals": current_llgs.model_dump(mode="json")},
        checkpoint_complete=False, completed_branch_count=done, total_branch_count=total)
    publish(pending=trace)

def publish(pending=None, result=None):
    if result is not None:
        state = result.model_dump(mode="json", exclude_defaults=True)
    else:
        state = dict(schema_version=RESULT_SCHEMA_VERSION, configuration=config, completed_iterations=len(traces),
                     warnings=list(warnings),
                     final_high_level_goals=current_hlgs.model_dump(mode="json"),
                     final_low_level_goals=current_llgs.model_dump(mode="json"),
                     iterations=[trace.model_dump(mode="json", exclude_defaults=True)
                                 for trace in traces + ([pending] if pending is not None else [])])
    if evaluation_output_directory is not None:
        atomic_write(Path(evaluation_output_directory) / f"{file_prefix}results.json", compact_cycle(state))
    if on_state is not None:
        on_state(state)

def finish(reason, converged):
    result = GlobalGoalCycleResult(
        schema_version=RESULT_SCHEMA_VERSION,
        converged=converged, stop_reason=reason, completed_iterations=len(traces),
        final_high_level_goals=current_hlgs, final_low_level_goals=current_llgs,
        iterations=traces, warnings=warnings, configuration=config)
    publish(result=result)
    return result
\end{lstlisting}

\begin{lstlisting}[style=llmprompt,language=Python,caption={Ciclo reale di interrogazione dei tre evaluator: ogni risposta globale viene salvata prima di passare al modello successivo.},label={lst:orchestrator-evaluator-sequence}]
for model in BOTTOM_UP_EVALUATOR_MODELS:
    global_record, votes = _critic1_vote(
        prompt, model, branches, analysis,
        temperature=evaluator_temperature, manage_model=True)
    global_records.append(global_record)
    if on_global_vote:
        on_global_vote(global_record)
    for branch in branches:
        if branch in votes:
            records[branch].append(votes[branch])
\end{lstlisting}

\begin{lstlisting}[style=llmprompt,language=Python,caption={Adattamento della chiamata alla reflection e registrazione dei candidati della riparazione.},label={lst:orchestrator-top-down-call}]
def _run_original_top_down(*, initial_feedback, records, stage, **kwargs):
    """Record all original reflection results for the transactional handoff."""
    kwargs["call_function"] = _adapt_bottom_up_feedback(
        kwargs["call_function"], initial_feedback)
    result, score, critique, attempts, history = generate_response_with_reflection(
        max_attempts=MAX_ATTEMPTS, **kwargs)
    records[stage] = {
        "score": score, "critique": critique, "attempts": attempts,
        "score_history": history,
        "candidate": _model_payload(result),
    }
    return result
\end{lstlisting}

\begin{lstlisting}[style=llmprompt,language=Python,caption={Gestore rientrante dei modelli locali: stesso modello nelle chiamate annidate e rilascio al termine.},label={lst:orchestrator-model-lease}]
@contextmanager
def loaded_local_model(model_id: str, *, context_length: int | None = None):
    """One GPU model per process; nested leases of the same model reuse it."""
    global _active_local_model
    with _model_lock:
        if _active_local_model is not None:
            if _active_local_model != model_id:
                raise RuntimeError(f"Cannot load {model_id} while {_active_local_model} has an active lease")
            if context_length is not None and _local_load_configs.get(model_id, {}).get("loaded_context_length", 0) < context_length:
                raise RuntimeError("Cannot enlarge context inside an active nested model lease")
            yield model_id
            return
        with _loaded_local_model_instance(model_id, context_length=context_length) as selected:
            _active_local_model = selected
            try:
                yield selected
            finally:
                _active_local_model = None
\end{lstlisting}

\begin{lstlisting}[style=llmprompt,language=Python,caption={Client del generatore Qwen: messaggi, schema richiesto e validazione del risultato.},label={lst:orchestrator-generator-client}]
def generate_response(
    prompt: str, sys_prompt: str, response_format: type[BaseModel] | None = None,
    *, stateless_retries: bool = False,
) -> str | BaseModel:
    """Use fresh messages per call; optionally repeat invalid JSON without history."""
    system_content = sys_prompt
    if response_format is not None:
        system_content += ("\nReturn only one JSON object matching this schema:\n"
                           + json.dumps(response_format.model_json_schema(), ensure_ascii=False))
    messages = [{"role": "system", "content": system_content}, {"role": "user", "content": prompt}]
    last_exception = None
    with loaded_local_model(QWEN_GENERATOR_MODEL):
        for _ in range(LM_STUDIO_GENERATOR_ATTEMPTS):
            try:
                response = _qwen_client().chat.completions.create(
                    model=QWEN_GENERATOR_MODEL, messages=messages, temperature=0,
                    max_tokens=QWEN_GENERATOR_MAX_TOKENS)
            except Exception as exc:
                raise RuntimeError(f"Qwen locale {QWEN_GENERATOR_MODEL!r} su {LM_STUDIO_BASE_URL}: {exc}") from exc
            choice = response.choices[0]
            content = choice.message.content
            if choice.finish_reason == "length":
                raise RuntimeError("Qwen reached its output-token budget; no incomplete result was accepted. Increase QWEN_GENERATOR_MAX_TOKENS and the loaded context.")
            if not content:
                raise RuntimeError("Local Qwen returned no final response; verify reasoning and output-token budgets.")
            if response_format is None:
                return content
            try:
                return response_format.model_validate(_extract_json(content))
            except ValueError as exc:
                last_exception = exc
                if not stateless_retries:
                    messages.extend([{"role": "assistant", "content": content},
                                     {"role": "user", "content": f"Correct the complete JSON to match the schema. Validation error: {exc}"}])
    raise ValueError(f"Local Qwen produced invalid structured output after {LM_STUDIO_GENERATOR_ATTEMPTS} attempts") from last_exception
\end{lstlisting}

\begin{lstlisting}[style=llmprompt,language=Python,caption={Client degli evaluator: gestione del modello e invio della richiesta testuale o con schema JSON.},label={lst:orchestrator-evaluator-client}]
def generate_evaluator_response(
    prompt: str,
    sys_prompt: str,
    *,
    temperature: float = 0,
    model: str | None = None,
    response_schema: dict | None = None,
) -> EvaluatorResponse:
    """Load the selected local model and make one request without chat history."""
    selected = model.strip() if model else _ministral_model()
    with loaded_local_model(selected):
        return _generate_evaluator_response(
            prompt, sys_prompt, temperature=temperature, model=selected,
            response_schema=response_schema)


def _generate_evaluator_response(
    prompt: str,
    sys_prompt: str,
    *,
    temperature: float = 0,
    model: str | None = None,
    response_schema: dict | None = None,
) -> EvaluatorResponse:
    """One stateless chat-completions request, textual or schema-constrained.

    The top-down critic uses text; bottom-up evaluators supply a JSON schema.
    No native chat endpoint or free-text fallback is used for structured calls.
    """
    selected_model = model.strip() if model else _ministral_model()
    if not selected_model:
        raise ValueError("The LM Studio evaluator model id cannot be empty.")
    messages = [
        {"role": "system", "content": sys_prompt},
        {"role": "user", "content": prompt},
    ]
    structured = response_schema is not None
    max_tokens = BOTTOM_UP_CRITIC_MAX_TOKENS if structured else MINISTRAL_CRITIC_MAX_TOKENS
    metadata = {
        "endpoint": "/v1/chat/completions", "structured_output": structured,
        "model": selected_model, "temperature": temperature,
        "max_output_tokens": max_tokens,
        "prompt_characters": sum(len(message["content"]) for message in messages),
        **_local_load_configs.get(selected_model, {}),
    }
    options = {}
    if structured:
        options["response_format"] = {"type": "json_schema", "json_schema": {
            "name": "documentary_review", "strict": True, "schema": response_schema}}
    try:
        response = _ministral_client().chat.completions.create(
            model=selected_model, messages=messages, temperature=temperature,
            max_tokens=max_tokens, **options)
    except Exception as exc:
        raise EvaluatorGenerationError(
            f"Local critic {selected_model!r} on {LM_STUDIO_BASE_URL}: {exc}",
            "", metadata) from exc
    choice = response.choices[0]
    content = choice.message.content or ""
    metadata.update(
        finish_reason=choice.finish_reason,
        usage=response.usage.model_dump(mode="json") if response.usage is not None else None)
    if choice.finish_reason != "stop" or not content.strip():
        raise EvaluatorGenerationError(
            f"Critic response incomplete or unavailable (finish_reason={choice.finish_reason!r}). "
            "Check output-token and loaded-context budgets; no partial response accepted.",
            content, metadata)
    return EvaluatorResponse(content, metadata)
\end{lstlisting}

La richiesta strutturata viene ulteriormente verificata da
\texttt{\_critic1\_vote} prima di partecipare al voto. La gestione del retry
globale e l'applicazione degli aggiornamenti rimangono separate dal client HTTP.
L'adattatore che inserisce il feedback soltanto al primo tentativo e la funzione
di riparazione della collezione sono riportati anche nei
Listing~\ref{lst:bottom-up-feedback-adapter} e~\ref{lst:collection-repair-code}.

\subsection{Fasi del bottom-up: definizioni con annotazioni}
\label{sec:annotated-bottom-up-phases}
Le funzioni seguenti sono estratte dai moduli correnti. I commenti
numerati spiegano input, chiamate e controlli nel punto in cui avvengono;
non sono istruzioni aggiunte ai prompt e non modificano il runtime.

\paragraph{Ricostruzione dei rami.}
Il caricamento esterno mantiene Qwen disponibile mentre si attraversano
i gruppi di figli. La funzione di ricostruzione del singolo gruppo
costruisce il prompt e il risultato riporta gli LLG usati come sorgente.
\begin{lstlisting}[style=llmprompt,language=Python,caption={Ricostruzione dei rami e provenienza degli HLG ricostruiti.},label={lst:annotated-reconstruction}]
def reconstruct_all_branches(
    goal_lists: dict[str, list[LowLevelGoal]] | HighLevelGoals,
    low_level_goals: LowLevelGoals | None = None,
):
    """Return reconstructions from existing top-down models or grouped LLGs."""
    # [01] Accetta anche le collezioni top-down e raggruppa i figli per padre.
    if isinstance(goal_lists, HighLevelGoals):
        if low_level_goals is None:
            raise ValueError("The top-down LLG collection is required")
        goal_lists = {
            bid: children for bid, parent in indexed_high_level_goals(goal_lists).items()
            if (children := low_level_goals_for_hlg(parent, low_level_goals))
        }
    # [02] Un insieme vuoto non richiede chiamate al modello.
    if not goal_lists:
        return {}
    # [03] Mantiene lo stesso modello caricato per tutti i gruppi.
    with loaded_local_model(QWEN_GENERATOR_MODEL):
        results = {}
        for branch_id, low_level_goals in goal_lists.items():
            # [04] La funzione del singolo ramo usa soltanto le descrizioni dei figli.
            output = reconstruct_high_level_goal(low_level_goals)
            # [05] Registra testo ricostruito e nomi dei figli sorgente.
            results[branch_id] = BottomUpHighLevelGoal(
                branch_id=branch_id,
                reconstructed_high_level_goal=output.reconstructed_high_level_goal,
                source_low_level_goal_ids=[goal.name for goal in low_level_goals])
        return results
\end{lstlisting}

\begin{lstlisting}[style=llmprompt,language=Python,caption={Embedding e cosine: controlli numerici prima delle etichette.},label={lst:annotated-numerical-primitives}]
@lru_cache(maxsize=4096)
def _encode_text(text: str) -> tuple[float, ...]:
    # [01] Un testo vuoto non riceve un vettore sostitutivo.
    if not text.strip():
        raise ValueError("Cannot embed empty goal text")
    model = _embedding_model()
    tokens = model.tokenizer(text, truncation=False, add_special_tokens=True)["input_ids"]
    # [02] Un testo troppo lungo genera un errore invece di essere troncato.
    if len(tokens) > model.max_seq_length:
        raise ValueError(f"Goal has {len(tokens)} tokens; embedding limit is {model.max_seq_length}")
    # [03] Codifica il testo con normalizzazione del vettore.
    vector = np.asarray(model.encode(text, normalize_embeddings=True), dtype=float)
    norm = np.linalg.norm(vector)
    if vector.ndim != 1 or not np.isfinite(vector).all() or not np.isfinite(norm) or norm == 0:
        raise ValueError("Invalid embedding vector")
    return tuple(vector / norm)

def _cosine(left, right) -> float:
    a, b = np.asarray(left, dtype=float), np.asarray(right, dtype=float)
    if (a.shape != b.shape or a.ndim != 1 or not np.isfinite(a).all()
            or not np.isfinite(b).all() or np.linalg.norm(a) == 0 or np.linalg.norm(b) == 0):
        raise ValueError("Invalid vectors for cosine similarity")
    # [01] Normalizza i vettori e limita il risultato in [-1, 1].
    return float(np.clip(np.dot(a / np.linalg.norm(a), b / np.linalg.norm(b)), -1, 1))
\end{lstlisting}

\begin{lstlisting}[style=llmprompt,language=Python,caption={Confronti HLG prime contro tutti gli HLG originali del progetto.},label={lst:annotated-similarity-pass}]
def _analyze_similarities_once(branches: dict[str, HighLevelGoal],
                         reconstructions: dict[str, BottomUpHighLevelGoal], *,
                         parent_threshold=None, duplicate_threshold=None) -> dict:
    """Compare each reconstructed HLG with every current original HLG."""
    config = similarity_configuration(parent_threshold=parent_threshold, duplicate_threshold=duplicate_threshold)
    threshold = config["similarity_threshold"]
    records, vectors, originals, reconstruction_errors, original_errors = {}, {}, {}, {}, {}
    ids = list(branches)
    # Every original HLG remains a column, including parents without children.
    # [01] Codifica le descrizioni originali e le ricostruzioni disponibili.
    for goal_id, parent in branches.items():
        try:
            originals[goal_id] = _encode_text(parent.description)
        except Exception as exc:
            original_errors[goal_id] = f"{type(exc).__name__}: {exc}"
        if goal_id in reconstructions:
            try:
                vectors[goal_id] = _encode_text(reconstructions[goal_id].reconstructed_high_level_goal)
            except Exception as exc:
                reconstruction_errors[goal_id] = f"{type(exc).__name__}: {exc}"

    pairs = []
    # [02] Ogni riga parte da una ricostruzione; le colonne sono HLG originali.
    for goal_id in ids:
        comparisons = []
        for other in ids:
            row = {"other_branch_id": other, "score": None}
            if goal_id not in reconstructions:
                row.update(status="NOT_RECONSTRUCTABLE")
            elif goal_id in reconstruction_errors or other in original_errors:
                row.update(status="ERROR", error=reconstruction_errors.get(goal_id) or original_errors[other])
            else:
                try:
                    row.update(status="VALID", score=_cosine(vectors[goal_id], originals[other]))
                except Exception as exc:
                    row.update(status="ERROR", error=f"{type(exc).__name__}: {exc}")
            comparisons.append(row)
        # [03] Il confronto diagonale identifica il rapporto con il proprio padre.
        own = comparisons[ids.index(goal_id)]
        record = {"status": own["status"], "score": own["score"],
                  "relationship": ("EQUIVALENT" if own["score"] >= threshold else "NOT_EQUIVALENT")
                                  if own["status"] == "VALID" else None,
                  "reconstructable": goal_id in reconstructions,
                  "duplicate_label": None, "duplicate_candidates": [],
                  "original_hlg_comparisons": comparisons,
                  "original_hlg_comparisons_complete": all(row["status"] == "VALID" for row in comparisons)}
        if own.get("error"):
            record["error"] = own["error"]
        elif own["status"] == "NOT_RECONSTRUCTABLE":
            record["error"] = "No reconstruction available for numerical comparison"
        records[goal_id] = record

    for goal_id in ids:
        record = records[goal_id]
        for row in record["original_hlg_comparisons"]:
            other = row["other_branch_id"]
            if other == goal_id:
                continue
            # [04] Soglia, disponibilita e stesso attore delimitano i candidati duplicati.
            candidate = (row["status"] == "VALID" and row["score"] >= threshold
                         and record["status"] == "VALID" and records[other]["status"] == "VALID"
                         and _actor_key(branches[goal_id]) == _actor_key(branches[other]))
            pairs.append({"branch_id": goal_id, **row, "candidate": candidate})
            if candidate:
                record["duplicate_candidates"].append(other)
        if record["status"] == "VALID":
            # [05] La candidatura numerica resta da confermare sulla documentazione.
            record["duplicate_label"] = "POTENTIAL_DUPLICATE" if record["duplicate_candidates"] else "NO_DUPLICATE_MATCH"
    config["parent_hlg_comparison_count"] = len(reconstructions)
    config["original_hlg_comparison_count"] = len(reconstructions) * len(branches)
    config["resolved_revision"] = _resolved_revision()
    return {"configuration": config, "branches": records, "pairs": pairs}
\end{lstlisting}

\begin{lstlisting}[style=llmprompt,language=Python,caption={Recupero degli errori tecnici della similarita, con budget separato.},label={lst:annotated-similarity-retry}]
def analyze_similarities(branches: dict[str, HighLevelGoal] | HighLevelGoals,
                         reconstructions: dict[str, BottomUpHighLevelGoal], *,
                         parent_threshold=None, duplicate_threshold=None) -> dict:
    """Recover failed numerical work before allowing documentary review.

    Successful text vectors stay cached. Errors are never cached by lru_cache;
    model/import failures can be retried after resetting their cached state.
    No text is shortened and no score or threshold is substituted.
    Missing reconstructions are unavailable inputs, not retryable numerical errors.
    """
    if isinstance(branches, HighLevelGoals):
        branches = indexed_high_level_goals(branches)
    history = []
    # [01] Riprova soltanto entro il budget tecnico della similarita.
    for attempt in range(1, SIMILARITY_TECHNICAL_ATTEMPTS + 1):
        analysis = _analyze_similarities_once(branches, reconstructions,
            parent_threshold=parent_threshold, duplicate_threshold=duplicate_threshold)
        errors = {}
        for bid, record in analysis["branches"].items():
            failed = [row for row in record.get("original_hlg_comparisons", [])
                      if row.get("status") == "ERROR"]
            if record["status"] == "ERROR" or failed:
                errors[bid] = {"parent_error": record.get("error"),
                               "comparison_errors": failed}
        history.append({"attempt": attempt, "accepted": not errors, "errors": errors})
        # [02] Senza errori tecnici termina i retry; un ramo puo restare non ricostruibile.
        if not errors:
            break
        print(f"[bottom-up] similarity recovery | attempt={attempt}/"
              f"{SIMILARITY_TECHNICAL_ATTEMPTS} | failed_branches={len(errors)}", flush=True)
        if attempt < SIMILARITY_TECHNICAL_ATTEMPTS:
            # [03] Ricarica il modello dopo errori; conserva i vettori validi in cache.
            _embedding_model.cache_clear()
            _runtime_import_error.cache_clear()
    # [04] Conserva la storia dei tentativi anche dopo il recupero.
    analysis["technical_attempts"] = history
    analysis["configuration"]["similarity_technical_attempts"] = SIMILARITY_TECHNICAL_ATTEMPTS
    return analysis
\end{lstlisting}

\begin{lstlisting}[style=llmprompt,language=Python,caption={Ingresso degli evaluator: testi, etichette e coppie candidate.},label={lst:annotated-review-payload}]
def build_review_payload(payload: dict) -> dict:
    """Expose parent labels and explicit candidate pairs; retain matrices in diagnostics."""
    branches = payload["branches"]
    # [01] Separa gli ID richiesti dalle informazioni di contesto della gerarchia.
    return {
        "requested_branch_ids": [g["branch_id"] for g in branches if g["status"] == "VALID"],
        "branches": [{"branch_id": g["branch_id"], "actor": g["actor"],
            "hlg": g["hlg"], "hlg_prime": g["hlg_prime"], "status": g["status"],
            "reconstructable": g["reconstructable"], "relationship": g["relationship"],
            "original_hlg_comparisons_complete": g["original_hlg_comparisons_complete"],
            "duplicate_candidates": g["duplicate_candidates"], "duplicate_label": g["duplicate_label"]}
            for g in branches],
        # [02] Ogni coppia contiene HLG prime e il suo specifico HLG candidato.
        "potential_duplicates": build_duplicate_candidate_pairs(branches),
        "duplicate_review": "explicit_prime_original_candidate_pairs_documentary_confirmation",
    }
\end{lstlisting}

\begin{lstlisting}[style=llmprompt,language=Python,caption={Risposta globale di un evaluator: chiamata, validazione e retry.},label={lst:annotated-model-vote}]
def _critic1_vote(prompt, model, branches, analysis, *, temperature=None, manage_model=False):
    """Share one per-iteration budget across load, context and output failures."""
    temperature = BOTTOM_UP_EVALUATOR_TEMPERATURE if temperature is None else temperature
    # [01] Richiede decisioni soltanto per confronti numerici disponibili.
    eligible = [b for b in branches if analysis["branches"][b]["status"] == "VALID"]
    expected = {b for b in eligible}
    schema = _documentary_response_schema(eligible, analysis)
    history, accepted, errors = [], {}, {}
    raw, payload, retry = None, None, prompt
    requested_context = LM_STUDIO_EVALUATOR_CONTEXT_LENGTH
    # [02] Lo stesso budget copre caricamento, contesto e risposta.
    for attempt in range(1, BOTTOM_UP_EVALUATOR_VOTE_ATTEMPTS + 1):
        if not eligible:
            break
        print(f"[bottom-up] global review | model={model} | attempt={attempt}/"
              f"{BOTTOM_UP_EVALUATOR_VOTE_ATTEMPTS}", flush=True)
        current_errors, generation, candidate = {}, {}, {}
        raw, payload, context_error = None, None, False
        stage = "loading"
        try:
            options = ({"context_length": requested_context}
                       if requested_context != LM_STUDIO_EVALUATOR_CONTEXT_LENGTH else {})
            lease = loaded_local_evaluator_model(model, **options) if manage_model else nullcontext()
            with lease:
                stage = "generation"
                # [03] Invia modello, temperatura, prompt e schema al client locale.
                raw = generate_evaluator_response(retry, SYSTEM_PROMPT,
                    temperature=temperature, model=model, response_schema=schema)
                generation = dict(getattr(raw, "metadata", {}))
                stage = "release"
            stage = "validation"
            # Also reject truncation if a custom client returns it instead of raising.
            if generation.get("finish_reason", "stop") != "stop":
                raise ValueError("Incomplete global response; regenerate complete concise JSON")
            payload = _json_object(raw)
            items = payload["decisions"]
            if any(not isinstance(item, dict) for item in items):
                raise ValueError("Every decision must be an object")
            counts = Counter(item.get("branch_id") for item in items
                             if isinstance(item.get("branch_id"), str))
            # [04] Una risposta deve coprire esattamente gli ID richiesti una sola volta.
            if set(counts) != expected or any(n != 1 for n in counts.values()) or len(items) != len(expected):
                raise ValueError("Return every requested branch exactly once; no extra or missing IDs")
            by_id = {item["branch_id"]: item for item in items}
            for branch in eligible:
                # [05] Controlla ciascuna azione prima di accettare la risposta intera.
                vote = _validate_decision(by_id[branch], branch, analysis)
                candidate[branch] = vote
        except Exception as exc:
            raw = getattr(exc, "raw_response", raw)
            generation = dict(getattr(exc, "metadata", generation))
            current_errors["response"] = f"{type(exc).__name__}: {exc}"
            context_error = _is_context_length_error(exc)
            # [06] Un errore di contesto puo richiedere il raddoppio del contesto.
            if context_error:
                generation.update(error_type="CONTEXT_LENGTH_EXCEEDED", recovery_required=True)
            elif stage in {"loading", "release"}:
                generation.update(error_type="MODEL_LOAD_FAILED")
        # [07] Nessun voto parziale di un tentativo scartato viene riutilizzato.
        success = not current_errors and len(candidate) == len(expected)
        if raw is not None:
            raw = str(raw)  # Persist plain text, not the metadata-bearing str subclass.
        generation.update(accepted=success, global_retry_policy="discard_entire_invalid_response",
                          recovery_requested_context_length=requested_context,
                          attempt_budget_scope="per_model_per_iteration")
        history.append(dict(attempt=attempt, errors=current_errors,
                            valid_branch_ids=list(candidate) if success else [],
                            generation=generation, requested_branch_ids=[b for b in eligible]))
        errors = current_errors
        if success:
            accepted = candidate
            break
        print(f"[bottom-up] model={model} | global response rejected: {current_errors}", flush=True)
        if context_error:
            requested_context = max(requested_context, generation.get("loaded_context_length", 0) or 0) * 2
            retry = prompt
            continue
        if stage in {"loading", "release"}:
            continue
        reason = current_errors.get("response", "Invalid global response")[:350]
        retry = (prompt + "\nCORRECTION: The previous global response was discarded. "
                 "Regenerate ALL requested decisions as one complete concise JSON. "
                 "Keep the explanation focused and the critique actionable. Error: " + reason)
    # [08] Restituisce voti accettati e cronologia del modello.
    global_record = EvaluatorVoteRecord(
        model=model, status="VALID" if accepted else "INVALID",
        decisions=list(accepted.values()), error=errors.get("response"), attempts=history)
    return global_record, accepted
\end{lstlisting}

\begin{lstlisting}[style=llmprompt,language=Python,caption={Giudizi dei tre modelli, quorum per ramo e verifica dei superstiti.},label={lst:annotated-global-voting}]
def evaluate_all_branches(
    project_description,
    branches,
    reconstructions,
    on_evaluation=None,
    *,
    actors=None,
    on_global_vote=None,
    similarity_analysis=None,
    cosine_payload=None,
    few_shot_example_count=None,
    evaluator_temperature=None,
):
    if isinstance(branches, HighLevelGoals):
        branches = indexed_high_level_goals(branches)
    if not branches:
        return {}
    analysis = similarity_analysis if similarity_analysis is not None else analyze_similarities(branches, reconstructions)
    if all(analysis["branches"][b]["status"] == "ERROR" for b in branches):
        results = {b: _inconclusive(
            b, "Similarity unavailable: " + analysis["branches"][b]["error"],
            status="SIMILARITY_UNAVAILABLE") for b in branches}
        if on_evaluation:
            on_evaluation(len(results), len(branches), results)
        return results
    actors = actors or Actors(actors=list({branches[b].actor.name: branches[b].actor for b in branches}.values()))
    # [01] Costruisce un solo prompt di snapshot condiviso dai tre modelli.
    prompt = _critic1_prompt(project_description, actors, branches, reconstructions, analysis, cosine_payload,
                            few_shot_example_count=few_shot_example_count)
    records = {b: [] for b in branches}
    global_records = []
    # [02] Interroga i modelli in sequenza senza mostrare i voti precedenti.
    for model in BOTTOM_UP_EVALUATOR_MODELS:
        global_record, votes = _critic1_vote(
            prompt, model, branches, analysis,
            temperature=evaluator_temperature, manage_model=True)
        global_records.append(global_record)
        if on_global_vote:
            on_global_vote(global_record)
        for branch in branches:
            if branch in votes:
                records[branch].append(votes[branch])
    results = {}
    for branch in branches:
        local = records[branch]
        # [03] La chiave conta azione e superstite, non la sola azione.
        counts = Counter(_operation_key(vote) for vote in local)
        match = analysis["branches"][branch]
        # [04] Senza due operazioni concordi il ramo resta inconcludente.
        if not counts or counts.most_common(1)[0][1] < EVALUATOR_VOTING_QUORUM:
            result = _inconclusive(branch, "No two valid documentary actions agree.")
            if any(attempt.get("generation", {}).get("error_type") == "CONTEXT_LENGTH_EXCEEDED"
                   for record in global_records for attempt in record.attempts):
                result = _inconclusive(branch,
                    "Evaluator context exceeded; insufficient agreeing valid votes. No automatic action.",
                    status="CONTEXT_LENGTH_EXCEEDED")
        else:
            key, _ = counts.most_common(1)[0]
            # [05] Unisce solo i feedback dei voti che sostengono la maggioranza.
            vote = _merge_agreeing_judgments(local, key)
            if match["status"] != "VALID":
                result = _inconclusive(branch, "Similarity comparison unavailable.")
            else:
                result = GlobalGoalEvaluationResult(
                    **vote.model_dump(), evaluation_status="DECIDED")
        results[branch] = result
        print(f"[bottom-up] ensemble | branch={branch} | cosine={match.get('score')} | "
              f"relationship={match.get('relationship')} | votes={dict(counts)} | "
              f"decision={(result.decision.value if result.decision else 'INCONCLUSIVE')}", flush=True)
    # Check all survivors against the original aggregate, not a mutated order-dependent graph.
    # [06] Controlla i superstiti sullo stesso aggregato, senza dipendenze di ordine.
    kept = {key for key, result in results.items() if result.decision == HighLevelGoalDecision.KEEP}
    for branch in branches:
        survivor = results[branch].surviving_branch_id
        # [07] Respinge una duplicazione se il superstite non e confermato KEEP.
        if survivor and survivor not in kept:
            results[branch] = _inconclusive(
                branch, "Duplicate survivor is not conclusively kept in this snapshot.",
                )
            print(f"[bottom-up] ensemble | branch={branch} | INCONCLUSIVE: "
                  "duplicate survivor is not kept", flush=True)
    if on_evaluation:
        on_evaluation(len(results), len(branches), results)
    return results
\end{lstlisting}

\begin{lstlisting}[style=llmprompt,language=Python,caption={Consolidamento del feedback dei soli voti concordi.},label={lst:annotated-feedback-merge}]
def _merge_agreeing_judgments(records, operation_key):
    """Join only the feedback prose of the agreeing evaluators for generation."""
    # [01] Scarta dal feedback i giudizi che non sostengono la stessa operazione.
    votes = [vote for vote in records if _operation_key(vote) == operation_key]
    if len(votes) < EVALUATOR_VOTING_QUORUM:
        raise ValueError("Merged feedback requires at least two agreeing valid votes")

    def join_text(parts):
        # [02] Elimina solo critiche testualmente identiche; non risolve contraddizioni.
        return "\n\n".join(dict.fromkeys(part.strip() for part in parts if part and part.strip()))

    updates = {}
    if votes[0].feedback:
        updates["feedback"] = {"critique": join_text(vote.feedback.critique for vote in votes)}
    return DocumentaryDecision.model_validate({**votes[0].model_dump(), **updates})
\end{lstlisting}

